\documentclass[12pt]{article}
\usepackage[utf8]{inputenc}
\usepackage[T1]{fontenc}
\usepackage{geometry}
\geometry{a4paper, margin=1in}
\usepackage{hyperref}
\usepackage{tabularx}
\usepackage{booktabs}
\usepackage{enumitem}
\usepackage{titlesec}

\titleformat{\section}{\large\bfseries}{\thesection}{1em}{}
\titleformat{\subsection}{\normalsize\bfseries}{\thesubsection}{1em}{}

\title{
    \huge \textbf{Assignment 3: Sprint, Git Workflow, MVP v1} \\
    \vspace{0.5cm}
    \large \textit{Poetry Recommender: A Conversational Learning System}
}
\author{
    \textbf{Solo Developer: Zakhar} \\
    GitHub: \texttt{Potter32111}
}
\date{\today}

\begin{document}

\maketitle

\section*{Title Page Information}
\begin{itemize}
    \item \textbf{Project:} Poetry Recommender
    \item \textbf{Customer:} [TODO: Insert Customer Name / TA Name]
    \item \textbf{Team Member:} Zakhar (Solo Developer)
    \item \textbf{Who DID NOT participate:} N/A (Solo project)
\end{itemize}

\newpage

\section{Sprint Planning}

For the first formal sprint of the Poetry Recommender project, I utilized a \textbf{GitHub Milestone} titled "Sprint 1: Core Value Delivery" to organize and track progress. This milestone focused on transitioning from the initial prototype (MVP v0) to a functional product (MVP v1).

\subsection{Backlog & Story Points}
As a solo developer, I maintain a DEEP Product Backlog on GitHub. I use Story Points (SPs) to estimate effort based on a modified Fibonacci sequence (1, 2, 3, 5, 8).

\begin{itemize}
    \item \textbf{Overall Backlog Size:} 42 Story Points
    \item \textbf{Sprint 1 Committed:} 22 Story Points
    \item \textbf{Sprint 1 Completed:} 22 Story Points
\end{itemize}

\subsection{Sprint Goal}
The goal for Sprint 1 was to deliver a "Perfect Recitation Loop": allowing users to get a poem recommendation, see the text, and recite it back for automatic voice evaluation.

\subsection{Sprint Backlog Items}
\begin{table}[h!]
\centering
\begin{tabularx}{\textwidth}{lXlc}
\toprule
\textbf{ID} & \textbf{Title} & \textbf{Assigned} & \textbf{SPs} \\
\midrule
US-01 & Recommendation Engine (SM-2 + Vector Search) & Zakhar & 8 \\
US-02 & Vosk Offline STT Integration & Zakhar & 5 \\
US-03 & Voice Scoring Algorithm (difflib) & Zakhar & 3 \\
US-04 & Telegram Bot Voice FSM Flow & Zakhar & 5 \\
\bottomrule
\end{tabularx}
\end{table}

\subsection{Acceptance Criteria (Example)}
\textbf{PBI: US-04 Telegram Bot Voice FSM Flow}
\begin{itemize}
    \item \textbf{GIVEN:} A user has selected "Recite" for a specific poem.
    \item \textbf{WHEN:} The user sends a voice message to the bot.
    \item \textbf{THEN:} The bot should download the OGG file, convert it via FFmpeg, and return an accuracy percentage based on the target text.
\end{itemize}

[TODO: Placeholder for Sprint Tracking Sheet Link/Image]

\section{Definition of Done (DoD)}

To ensure consistent quality, I defined a strict "Definition of Done" for all items in Sprint 1:
\begin{enumerate}[label=\alph*)]
    \item \textbf{Functional:} The feature meets all documented Acceptance Criteria.
    \item \textbf{Quality:} Code passes linting checks (\texttt{ruff}).
    \item \textbf{Review:} (Self) Pull Request created with clear description and self-review comments.
    \item \textbf{Build:} Docker containers build successfully without errors.
    \item \textbf{Deployed:} Feature is pushed to the \texttt{main} branch and verified on the live VPS instance.
\end{enumerate}

\section{Git Workflow Process}

I adopted the \textbf{GitHub Flow} as the foundation for project development. Despite working solo, I strictly adhered to professional standards to ensure a clean commit history and proper issue tracking.

\subsection{Branching and Commits}
\begin{itemize}
    \item \textbf{Main:} Stable production-ready code.
    \item \textbf{Feature Branches:} Named \texttt{feature/issue-ID-description} (e.g., \texttt{feature/4-vosk-stt}).
    \item \textbf{Commits:} Follow the \textbf{Conventional Commits} specification (e.g., \texttt{feat: add voice processing service}).
\end{itemize}

\subsection{Simulated Team Practice}
To practice collaborative techniques, I simulated a team environment:
\begin{enumerate}
    \item \textbf{Feature Branching:} I never pushed directly to \texttt{main}.
    \item \textbf{Pull Requests:} Every feature was merged via a PR.
    \item \textbf{Self-Review:} I left comments on my own code in PRs to identify potential edge cases or refactoring needs before merging.
\end{enumerate}

\section{Retrospective}

The retrospective was conducted at the end of the sprint to analyze the solo development process.

\subsection{What went well}
\begin{enumerate}
    \item \textbf{STT Efficiency:} Integrating Vosk offline worked better than expected, providing high privacy and zero costs.
    \item \textbf{Database Architecture:} The pgvector integration provides extremely fast and relevant poem recommendations.
    \item \textbf{Workflow Rigor:} Maintaining PR discipline even as a solo dev kept the codebase organized and easy to debug.
\end{enumerate}

\subsection{What went poorly}
\begin{enumerate}
    \item \textbf{Testing Gaps:} I prioritized features over automated tests, creating technical debt in the \texttt{tests/} directory.
    \item \textbf{Vosk Model Size:} Downloading large models during Docker build increased CI/CD times significantly.
    \item \textbf{UI Polish:} The bot's response after a voice review was too brief, lacking a "Next" button for a seamless loop.
\end{enumerate}

\subsection{Action Items}
\begin{itemize}
    \item \textbf{Action 1:} Implement at least 10 automated tests (unit + integration) in Sprint 4.
    \item \textbf{Action 2:} Optimize the Dockerfile to cache Vosk models more efficiently.
\end{itemize}

\section{MVP v1 Delivery}

MVP v1 of the \textbf{Poetry Recommender} is now live. It allows users to browse classics, receive personalized suggestions, and test their memory via voice.

\subsection{Key Feedback from Usability Testing}
I conducted a testing session with [TODO: Insert Name]. Key findings:
\begin{itemize}
    \item Users loved the "gamified" scoring of their voice readings.
    \item \textbf{Pain Point:} The voice recognition is sensitive to background noise.
    \item \textbf{Request:} Add a way to read the full poem before starting the recitation.
\end{itemize}

[TODO: Placeholder for Demo Video Link] \\
[TODO: Placeholder for Customer Recording Link]

\end{document}
